\section*{Problem 1}

在 $\mathbb{R}^2$ 上定义无穷范数
\[
\|(x,y)\|_\infty = \max\{|x|,|y|\}.
\]
设映射 $\mathcal{T} : \mathbb{R}^2 \to \mathbb{R}^2$ 定义为
\[
\mathcal{T} \left(x,y\right) = \left(\frac14 x + \frac12 y + 1,\; \frac13 x + \frac16 y - 2\right)^\top.
\]
证明 $\mathcal{T}$ 是压缩映射，并求其唯一不动点。

\begin{proof}
在标准基下，$\mathcal{T}$ 可表示为
\[
\mathcal{T} \left(\bvec{x}\right) = A \bvec{x} + \bvec{b}
\]
其中
\[A = \begin{bmatrix}
\frac14 & \frac12 \\
\frac13 & \frac16
\end{bmatrix}, \quad \bvec{b} = \begin{pmatrix}
1 \\
-2
\end{pmatrix}.\]
对任意 $\bvec{x} = (x_1, x_2)^\top, \bvec{y} = (y_1, y_2)^\top \in \mathbb{R}^2$，有
\[\begin{aligned}
\|\mathcal{T}(\bvec{x}) - \mathcal{T}(\bvec{y})\|_\infty &= \|A(\bvec{x} - \bvec{y})\|_\infty \\
&= \max\left\{\left|\frac14 (x_1 - y_1) + \frac12 (x_2 - y_2)\right|, \left|\frac13 (x_1 - y_1) + \frac16 (x_2 - y_2)\right|\right\} \\
&\leq \max\left\{\frac14 |x_1 - y_1| + \frac12 |x_2 - y_2|, \frac13 |x_1 - y_1| + \frac16 |x_2 - y_2|\right\} \\
&\leq \max\left\{\frac34 \max\{|x_1 - y_1|, |x_2 - y_2|\}, \frac12 \max\{|x_1 - y_1|, |x_2 - y_2|\}\right\} \\
&= \frac34 \max\{|x_1 - y_1|, |x_2 - y_2|\} \\
&= \frac34 \|\bvec{x} - \bvec{y}\|_\infty.
\end{aligned}\]
因此 $\mathcal{T}$ 是压缩映射。

为了求唯一不动点，我们解方程 $\mathcal{T}(\bvec{x}) = \bvec{x}$，即
\[
\begin{bmatrix}
\frac14 & \frac12 \\
\frac13 & \frac16
\end{bmatrix}
\begin{pmatrix}
x_1 \\
x_2
\end{pmatrix} + \begin{pmatrix}
1 \\
-2
\end{pmatrix} = \begin{pmatrix}
x_1 \\
x_2
\end{pmatrix}.
\]
这等价于
\[
\begin{cases}
\frac14 x_1 + \frac12 x_2 + 1 = x_1, \\
\frac13 x_1 + \frac16 x_2 - 2 = x_2.
\end{cases}
\]
解得 $x_1 = -\dfrac{4}{11}, x_2 = -\dfrac{28}{11}$。因此唯一不动点为 $\boxed{\left(-\frac{4}{11}, -\frac{28}{11}\right)^\top}$。

\end{proof}

\section*{Problem 2}

设 $\mathcal{T} : \mathbb{R} \to \mathbb{R}$ 定义为
\[
\mathcal{T}(x) = \frac12 \cos x.
\]
证明方程
\[
x = \frac12 \cos x
\]
有唯一实根，并说明迭代
\[
x_{n+1} = \mathcal{T}(x_n)
\]
对任意初值 $x_0 \in \mathbb{R}$ 都收敛到该实根。

\begin{proof}
首先证明 $\mathcal{T}$ 是压缩映射。对任意 $x, y \in \mathbb{R}$，有
\[
\begin{aligned}
|\mathcal{T}(x) - \mathcal{T}(y)| &= \left|\frac12 \cos x - \frac12 \cos y\right| = \frac12 |\cos x - \cos y| \\
&= \frac12 \left| {\left(\cos \frac{x + y}{2} \cos \frac{x - y}{2} - \sin \frac{x + y}{2} \sin \frac{x - y}{2}\right)- \left(\cos \frac{x + y}{2} \cos \frac{x - y}{2} + \sin \frac{x + y}{2} \sin \frac{x - y}{2}\right)}\right| \\
&= \frac12 \left| -2 \sin \frac{x + y}{2} \sin \frac{x - y}{2} \right| = \left|\sin \frac{x + y}{2}\right| \cdot \left|\sin \frac{x - y}{2}\right| \leq \left|\sin \frac{x - y}{2}\right| \leq \frac12 \left|x - y\right|.
\end{aligned}
\]
因此 $\mathcal{T}$ 是压缩映射。对任意初值 $x_0 \in \mathbb{R}$，迭代得序列
\[
x_1 = \mathcal{T}(x_0),\quad x_2 = \mathcal{T}(x_1),\quad \ldots, \quad x_n = \mathcal{T}(x_{n-1}).
\]
我们断言此序列收敛到 $\mathcal{T}$ 的唯一不动点 $x^*$，即上述序列有极限
\[
\lim_{n \to \infty} x_n = x^*
\]
且满足 $x^* = \mathcal{T}(x^*)$。先证明存在性，记 $\mu \left(x\right) = \mathcal{T} \left(x\right) - x$，代入有
\[
\mu \left(0\right) = \frac{1}{2} > 0, \quad \mu \left(1\right) = \frac{1}{2} \cos 1 - 1 < -\frac{1}{2} < 0.
\]
由零点定理见 $x^* \in \left(0, 1\right)$。

再证明唯一性，设 $x_1^*, x_2^*$ 是 $\mathcal{T}$ 的两个不动点，则
\[\begin{aligned}
|x_1^* - x_2^*| = |\mathcal{T}(x_1^*) - \mathcal{T}(x_2^*)| \leq \frac12 |x_1^* - x_2^*|.
\end{aligned}\]
移项得 $\frac12|x_1^*-x_2^*|\leq 0$，故 $x_1^* = x_2^*$。

最后证明收敛性，利用 $\mathcal{T}$ 的压缩性质有
\[
\begin{aligned}
    \left|x_n - x^*\right| = \left| \mathcal{T} \left(x_{n - 1}\right) - \mathcal{T} \left(x^*\right) \right| &\leq \frac12 \left|x_{n - 1} - x^*\right| \\
    &\leq \frac12 \cdot \frac12 \left|x_{n - 2} - x^*\right| = \left(\frac12\right)^2 \left|x_{n - 2} - x^*\right| \\
    &\leq \cdots \leq \left(\frac12\right)^{n} \left|x_0 - x^*\right|.
\end{aligned}
\]
由于 $\left(\frac12\right)^n|x_0-x^*|\to0$，故 $\lim_{n \to \infty} x_n = x^*$。因此方程 $x = \frac12 \cos x$ 有唯一实根，并且对任意初值 $x_0 \in \mathbb{R}$，迭代 $x_{n+1} = \mathcal{T}(x_n)$ 都收敛到该实根。


\end{proof}
\section*{Problem 3}

设
\[
A =
\begin{bmatrix}
1 & 2 \\
2 & 4
\end{bmatrix}.
\]
在 $\mathbb{R}^2$ 中取 Euclid 范数 $\|\cdot\|_2$，求算子范数 $\|A\|_2$。

\begin{proof}
算子范数 $\|A\|_2$ 定义为
\[\|A\|_2 = \sup_{\bvec{x} \neq \bvec{0}} \frac{\|A\bvec{x}\|_2}{\|\bvec{x}\|_2}.\]
由于 $A$ 是对称矩阵，$\|A\|_2$ 等于 $A$ 的最大特征值的绝对值。计算 $A$ 的特征值有
\[\det \begin{bmatrix}
1 - \lambda & 2 \\
2 & 4 - \lambda
\end{bmatrix} = (1 - \lambda)(4 - \lambda) - 4 = \lambda^2 - 5\lambda = 0.\]
因此 $A$ 的特征值为 $\lambda_1 = 0, \lambda_2 = 5$。最大特征值的绝对值为 $5$，因此 $\|A\|_2 = \boxed{5}$。


\end{proof}
\section*{Problem 4}

考虑初值问题
\[
y^\prime \left(x\right) = y\left(x\right),\quad y\left(0\right) = 1.
\]
将其改写为积分方程
\[
y\left(x\right) = 1 + \int_0^x y\left(t\right)\, \mathrm{d} t.
\]
在区间 $[0,a]$ 上的连续函数空间 $C\left[0,a\right]$ 中取范数
\[
\left\|y\right\|_\infty = \max_{x \in \left[0,a\right]} \left|y\left(x\right)\right|.
\]
定义算子
\[
\left(Ay\right)\left(x\right) = 1 + \int_0^x y\left(t\right)\, \mathrm{d} t.
\]
证明当 $0 < a < 1$ 时，$A$ 为压缩映射，从而该初值问题在 $C\left[0,a\right]$ 中有唯一解。

\begin{proof}
若 $y\in C[0,a]$，则 $Ay$ 仍为连续函数，因此 $A$ 是 $C[0,a]$ 到自身的映射。又 $C[0,a]$ 在无穷范数下是完备空间。
任取 $y_1 \left(x\right), y_2 \left(x\right) \in C\left[0,a\right]$，有
\[
\left(Ay_1\right)\left(x\right) - \left(Ay_2\right)\left(x\right) = \int_0^x \left(y_1\left(t\right) - y_2\left(t\right)\right)\, \mathrm{d} t.
\]
因此
\[
\left\|\left(Ay_1\right) - \left(Ay_2\right)\right\|_\infty = \max_{x \in \left[0,a\right]} \left|\int_0^x \left(y_1\left(t\right) - y_2\left(t\right)\right)\, \mathrm{d} t\right|.
\]
由积分不等式得
\[
\left|\int_0^x \left(y_1\left(t\right) - y_2\left(t\right)\right)\, \mathrm{d} t\right| \leq \int_0^x \left|y_1\left(t\right) - y_2\left(t\right)\right|\, \mathrm{d} t \leq \int_0^x \left\|y_1 - y_2\right\|_\infty \, \mathrm{d} t = a \left\|y_1 - y_2\right\|_\infty.
\]
所以
\[
\left\|\left(Ay_1\right) - \left(Ay_2\right)\right\|_\infty \leq a \left\|y_1 - y_2\right\|_\infty.
\]
当 $0 < a < 1$ 时，$A$ 为压缩映射。由 Banach 不动点定理，$A$ 在 $C[0,a]$ 中有唯一不动点，从而该初值问题有唯一解。
\end{proof}
